\section{Introduction}\label{sec:introduction}

A pure type system specifies which sorts classify other sorts and which
sorts may form dependent function types.  A term is \emph{weakly
$\beta$-normalizing} if some sequence of $\beta$-reductions reaches a normal
form; it is \emph{strongly $\beta$-normalizing} if every such sequence is
finite.  The distinction is between existence of a terminating computation
and termination under every choice of reductions.

We use the full annotated syntax: reduction is allowed in the domain of
an abstraction and in both components of a dependent product. Contexts may
be open. A legal expression is either side of a derivable typing judgment;
declaration types are included by context validity. The formal rules and
conventions are given below before the main theorem.

\subsection{Pure type systems and normalization}\label{pts:setup}

We first fix the syntax and the scope of the normalization properties.
All contexts in the typing rules are finite, and all reductions include
reductions inside annotations.

\begin{definition}[Specification and expressions]\label{pts:specification}
A \emph{pure type specification} is a triple
$\mathcal P=(\mathcal S,\mathcal A,\mathcal R)$, where $\mathcal S$ is a
set of \emph{sorts}, $\mathcal A\subseteq\mathcal S^2$ is a relation of
\emph{axioms}, and $\mathcal R\subseteq\mathcal S^3$ is a relation of
\emph{product rules}. Neither relation is assumed functional.
Using a countably infinite set of variables disjoint from $\mathcal S$, form the
raw expressions
\[
 M,N,A,B ::= x\mid s\mid MN\mid\lambda x:A.M\mid\Pi x:A.B,
 \qquad s\in\mathcal S.
\]
Expressions are identified up to renaming bound variables. In each binder
the variable binds in the body or codomain, but not in its annotation or
domain. Application associates to the left. Substitution is capture
avoiding; $M[x:=N]$ denotes substitution for one free variable.

A raw context is a finite sequence
$\Gamma=(x_1:A_1,\ldots,x_n:A_n)$ with distinct declared variables.
Its domain is $\operatorname{dom}(\Gamma)=\{x_1,\ldots,x_n\}$.
Whenever a rule appends $x:A$, its variable $x$ is fresh for the preceding
context and for $A$.
\end{definition}

\begin{definition}[Full beta reduction]\label{pts:beta}
The relation $\to_\beta$ is the compatible closure of
\[
             (\lambda x:A.M)N\to_\beta M[x:=N].
\]
Compatibility permits a step in either child of an application, product,
or lambda. In particular, both
$\lambda x:A.M\to_\beta\lambda x:A'.M$ when $A\to_\beta A'$ and
$\Pi x:A.B\to_\beta\Pi x:A'.B$ are permitted.
Write $\to_\beta^*$ for its reflexive-transitive closure and
$=_\beta$ for the equivalence relation it generates. A normal expression
has no beta redex at any position. An expression is \emph{weakly
normalizing} if it reduces to a normal expression, and \emph{strongly
normalizing} if it starts no infinite reduction sequence.
\end{definition}

Typing is the least relation closed under these seven rules, with
$s,s_1,s_2,s_3$ literal sorts.
\begin{gather*}
 \frac{(s_1,s_2)\in\mathcal A}{\varnothing\vdash s_1:s_2}
 \;\textsc{Axiom}
 \qquad
 \frac{\Gamma\vdash A:s}{\Gamma,x:A\vdash x:A}
 \;\textsc{Variable}
 \\
 \frac{\Gamma\vdash M:B\qquad\Gamma\vdash A:s}
      {\Gamma,x:A\vdash M:B}
 \;\textsc{Weakening}
 \\
 \frac{\Gamma\vdash A:s_1\qquad
       \Gamma,x:A\vdash B:s_2\qquad(s_1,s_2,s_3)\in\mathcal R}
      {\Gamma\vdash\Pi x:A.B:s_3}
 \;\textsc{Product}
 \\
 \frac{\Gamma,x:A\vdash M:B\qquad\Gamma\vdash\Pi x:A.B:s}
      {\Gamma\vdash\lambda x:A.M:\Pi x:A.B}
 \;\textsc{Abstraction}
 \\
 \frac{\Gamma\vdash M:\Pi x:A.B\qquad\Gamma\vdash N:A}
      {\Gamma\vdash MN:B[x:=N]}
 \;\textsc{Application}
 \\
 \frac{\Gamma\vdash M:A\qquad\Gamma\vdash B:s\qquad A=_\beta B}
      {\Gamma\vdash M:B}
 \;\textsc{Conversion}.
\end{gather*}
The relation is understood relative to the fixed specification.
When two specifications are compared we indicate the system on the
judgment. A derivation is always finite, including a finite conversion
chain witnessing each use of $=_\beta$.

\begin{definition}[Contexts, legality, and system normalization]
\label{pts:normalization}
A context $(x_1:A_1,\ldots,x_n:A_n)$ is \emph{valid} if each $A_i$ is
typed at some literal sort in the preceding prefix. An expression is
\emph{sorted} in $\Gamma$ if $\Gamma\vdash A:s$ for some $s\in\mathcal S$.
It is \emph{legal} in $\Gamma$ if it is the subject or the expected type
of a derivable judgment in $\Gamma$.

The system is weakly normalizing, respectively strongly normalizing, if
every legal expression in every context is weakly normalizing,
respectively strongly normalizing. These are system-wide properties;
neither statement replaces its universal quantifier by a hypothesis
about a single expression. The reduction relation acts on the expression
itself, including its annotations; it does not unfold declarations in
the external context.
\end{definition}



\begin{theorem}[The $\beta$-Barendregt--Geuvers--Klop conjecture]
\label{thm:main}
Let $\mathcal P=(\mathcal S,\mathcal A,\mathcal R)$ be any pure type system.
If every legal expression of $\mathcal P$, in every valid context, is weakly
$\beta$-normalizing, then every such expression is strongly
$\beta$-normalizing. No functionality assumption is imposed on
$\mathcal A$ or $\mathcal R$.
\end{theorem}

Both quantifiers in Theorem~\ref{thm:main} range over the whole system.
The theorem does not assert that an individual term with a normal form is
strongly normalizing. In particular, the proof may use weak normalization
of expressions other than the one whose strong normalization is being
established. The result concerns $\beta$-reduction; it makes no assertion
about adding $\eta$-reduction.

\paragraph{History and significance.}
Geuvers formulates the system-wide conjecture, for both $\beta$ and
$\beta\eta$ reduction, in his thesis \cite[Conjecture~8.1.2]{Geuvers1993}.
Theorem~\ref{thm:main} establishes its $\beta$ instance for arbitrary
specifications. It allows a proof that every legal expression has a normal
form to serve also as a proof of termination under every choice of beta
reductions. This includes computations inside types and annotations,
which may continue after the computational body has reached normal form.
The hypothesis remains a property of the whole system: the argument can
use normal forms of open expressions other than the expression under study.

One line of earlier work reduces strong normalization to weak
normalization of translated terms. The translation must preserve typing,
so the products available in the source specification are decisive.
S{\o}rensen's continuation-passing construction proves the implication for
generalized nondependent, clean, negatable systems
\cite[Theorem~3.5.20]{Sorensen1997}; the uniform nondependent result of
Barthe, Hatcliff and S{\o}rensen develops this approach
\cite{BartheHatcliffSorensen2001}. Mull's thesis weakens the cleanliness
and negatability restrictions within the nondependent tiered setting
\cite[Theorem~2]{MullThesis2023}. Here the sorts form a finite axiom chain
and each product inherits its codomain's sort. In particular, the theorem
allows a weakly clean alternative at each non-top tier that does not require
negatability, alongside an alternative retaining negatability and a
further cleanliness condition. These results explain both the force of
type-preserving translations and the importance of the product rules they
require.

Structural transformations offer a complementary way to reduce the
problem. Roux and van Doorn prove weak-normalization preservation for
disjoint sums and a specified family of added product rules, using
labelled syntax and dependency erasure
\cite[Sections~3--4]{RouxVanDoorn2014}. Mull proves that the
Barendregt--Geuvers--Klop implication transfers from the irrelevance
reduction of a tiered system to that system
\cite[Theorem~47]{Mull2023}. This is a conditional reduction of the
conjecture within the tiered class. More recently, Roux proposes an
internal proof translation with reported constructions for the simply
typed lambda calculus, System~F, and a trivial recursive type
\cite[slides~16--17]{Roux2025}. The general specification considered here
need not have a hierarchy of sorts, functional axiom or product relations,
or the products required by any of these translations.

Annotations impose a separate constraint on comparisons between systems.
Barthe and Coquand exhibit nonnormalizing pure type systems in which
erasing lambda-domain annotations identifies terms of the same type and
context that are not beta-convertible
\cite[Theorem~17]{BartheCoquand2006}. The auxiliary erasure used below
retains all annotations. It removes only proof/data mode labels, and its
reduction-lifting property is proved for each ordinary typed beta step.

\paragraph{Proof strategy.}
A finite derivation of a counterexample uses only a finite part of the
specification, so the proof first reduces to finite sort and rule sets.
Under weak normalization and confluence, every legal expression then has
a unique normal form. The \emph{sort profile} of a normal type records
all sorts at which it can be typed. Retaining this set accommodates
nonfunctional specifications; tracking its growth under substitution
replaces an appeal to uniqueness of types.

The candidate argument follows the Tait--Girard reducibility tradition
\cite[Chapters~6 and~14]{Girard1989}. Its candidates are sets of strongly
normalizing terms closed under the head expansions needed here; their
lattice, product tests and substitution laws are proved explicitly,
without requiring reduction closure. Profiles form a finite signed graph:
in a dependent product, the domain reverses inclusion between candidates
and the codomain preserves it. Processing its strongly connected
components in dependency order establishes that a strongly normalizing
function applied to any finite, correctly typed stack of strongly
normalizing arguments remains strongly normalizing. This excludes a
smallest nonnormalizing application and yields the theorem.

Components whose signs are consistent admit an ordered fixed point by
Tarski's theorem \cite{Tarski1955}. A second construction handles products
with at most one child in the current component. The remaining components
require observation tables, indexed by stacks of normal arguments and
auxiliary values. Negative domain occurrences make the equations for
these spaces circular. The construction resolves that circularity using
well-founded recursion for some profiles and alternating finite stages
for the others. Its key bound is uniform over all observations of a fixed
typed expression, although the possible arguments and stack lengths are
unbounded. Semantic substitution and reduction transport then give the
candidate equality required at a dependent application.

This last construction is conditional on a graph exclusion. The forbidden
configuration records the profile pattern forced if substitution exposes
new product nodes counted by the recursion's measure on normal types.
Weak normalization rules it out through a second, typed argument. The
product rules in such a configuration support a small
formula calculus and a relational realization of Hurkens's
well-foundedness paradox \cite{Hurkens1995}, using the presentation of
Geuvers \cite[Section~2]{Geuvers2007}.

Realizing that argument requires more than reproducing its final diagonal
steps. The derived formula calculus can quantify only over data types
supported by particular product rules in the profile graph. Moreover, the
general cancellation law for wrappers compares formula observations by
logical equivalence; this gives
no permission to replace a data expression inside a function.
Two data channels provide the predicates and observations needed by the
diagonal. Their access construction retains an open parameter exactly
while rounding a second data argument. Explicit validity and
extensionality conditions then state which observations and predicates
respect the resulting comparisons. The guarded diagonal uses precisely
these conditions, and every logical quantifier is checked against the
available product rules.

The resulting labelled term proves a designated terminal formula in a
context where a direct syntactic descent excludes every proof-mode term
with normal erasure. Data inhabitants of the same terminal do not affect
this statement. Erasure leaves all type annotations intact, and every
ordinary reduction of the erased typed proof lifts to a labelled
reduction. Weak normalization would therefore give the excluded normal
erasure. This contradiction supplies the graph exclusion and completes
the candidate argument.

\paragraph{Organization.}
Section~\ref{pts:metatheory} establishes the required PTS metatheory.
Sections~\ref{sec:profiles} and~\ref{sec:simple-components} develop the
candidate interface and its two simpler realizations.
Section~\ref{obs:section} constructs the remaining interpretation,
conditional on the graph exclusion. Sections~\ref{sec:enc-wrappers},
\ref{chan-section} and~\ref{rel-section} build the logical and relational
realization. Section~\ref{rel-diagonal} gives the guarded diagonal
argument. Section~\ref{sec:conclusion} proves the exclusion and completes
the theorem.
